\section{Conclusion}
\label{sec:conclusion}

A strategy can be redundant for current operation while remaining useful as a
carrier of a future research capability. A manager who wants to reduce active
maintenance burden must specify both roles before deleting entries. Under the
declared finite factorization, preserving the operating frontier and generative
closure preserves the library's productive process. The theorem's assumptions
are substantive: dependence on a named strategy or another omitted attribute
can invalidate that reduction.

Under identity closure, minimum-burden safe retention is a tagged weighted
cover. This gives access to established exact and approximate methods, with
independent checks in the original strategy coordinates. Safe local deletion
alone provides no general burden guarantee, as the heaviest-first family
shows. The registered benchmark confirms exact-method agreement where
available and documents burden quality and preprocessing effects within its
realized design; it does not support general speed rankings.

The point-in-time application preserves many research options at lower declared
burden, yet the protected common-equity policy makes no additional investment
and its one ETF replacement loses held-out value. The synthetic experiment
shows that small positive bridge value can also be difficult to detect under
the specified noise. Preserving an option therefore leaves a separate
statistical and economic decision about exercise.

The natural extensions are capability interactions beyond identity closure and
an independent financial study with documented institutional capabilities and
maintenance costs measured before future outcomes. The present findings support
a precise retention contract, rather than a claim of alpha or deployable
performance.
